\documentclass[12pt]{article}

\usepackage{amsmath, amssymb, amsthm, mathtools}
\usepackage{hyperref}

\newtheorem{theorem}{Theorem}
\newtheorem{lemma}[theorem]{Lemma}
\newtheorem{proposition}[theorem]{Proposition}
\newtheorem{corollary}[theorem]{Corollary}
\newtheorem{definition}[theorem]{Definition}

\newcommand{\VV}{\mathcal{V}}
\newcommand{\EE}{\mathcal{E}}
\newcommand{\GG}{\Gamma}


\begin{document}
\title{Submission D solution to Problem 10}
\date{}
\maketitle


\section*{Problem Statement}

Prove the following statement:

Let $\Gamma$ be a finite simplicial graph with $|\VV(\GG)|\geq 3$, and $(M_v,\varphi_v)$ be a separable tracial von Neumann algebra for each $v\in \VV$. If $\Gamma$ is irreducible (in the sense that there does not exists subgraphs $\GG_1,\GG_2$ such that $ \VV(\GG_1)\cup \VV(\GG_2)=\VV(\GG) $ and every $(v,u)\in \EE(\GG) $ for all $v\in\VV(\GG_1),u\in \VV(\GG_2)$), and each $M_v$ has a trace zero unitary, then the graph product von Neumann algebra $M$ is properly proximal.

\section*{Solution}

\begin{theorem}
Let $\Gamma$ be a finite simplicial graph with $|\VV(\Gamma)|\geq 3$, and $(M_v,\varphi_v)$ be a separable tracial von Neumann algebra for each $v\in \VV(\Gamma)$. If $\Gamma$ is irreducible (its complement graph $\Gamma^c$ is connected), and each $M_v$ has a trace zero unitary, then the graph product von Neumann algebra $M_\Gamma = \overline{\bigotimes}_\Gamma M_v$ is properly proximal.
\end{theorem}

\begin{proof}
We proceed by mathematical induction on the number of vertices $n = |\VV(\Gamma)|$. 

\medskip\noindent\textbf{Step 1: Base Cases ($n=3$)}

Let $\Gamma$ be an irreducible graph on 3 vertices. Irreducibility dictates that the complement graph $\Gamma^c$ is connected. The only connected graphs on 3 vertices are the complete graph $K_3$ and the path $P_3$. 
We rely on the foundational dimension criteria for proper proximality of free products (Boutonnet, Ioana, and Peterson, 2021, Theorem 4.1): The free product of two tracial von Neumann algebras $A \ast B$ is properly proximal provided $\dim A \ge 2$ and $\dim B \ge 3$. By hypothesis, each vertex algebra $M_v$ possesses a trace-zero unitary, guaranteeing $\dim M_v \ge 2$.
\begin{itemize}
    \item \textbf{Case 1a ($\Gamma^c = K_3$):} The graph $\Gamma$ is totally disconnected. Thus, $M_\Gamma = M_1 \ast M_2 \ast M_3$. We rewrite this via associativity as the free product of $A = M_1$ and $B = M_2 \ast M_3$. Since $\dim M_2 \ge 2$ and $\dim M_3 \ge 2$, their free product contains a diffuse subalgebra (isomorphic to the group von Neumann algebra of $\mathbb{Z}_2 \ast \mathbb{Z}_2$), yielding an infinite linear dimension, $\dim B = \infty \ge 3$. Since $\dim A \ge 2$, the free product $A \ast B$ strictly satisfies the threshold and is properly proximal.
    
    \item \textbf{Case 1b ($\Gamma^c = P_3$):} Let the complement edges be $(1,2)$ and $(2,3)$. The graph $\Gamma$ strictly contains only the single edge $(1,3)$. Thus, the graph product evaluates as $M_\Gamma = (M_1 \overline{\otimes} M_3) \ast M_2$. We set $A = M_1 \overline{\otimes} M_3$ and $B = M_2$. The tensor product yields $\dim A = \dim M_1 \times \dim M_3 \ge 2 \times 2 = 4 \ge 3$. Since $\dim B \ge 2$, the free product $A \ast B$ strictly satisfies the threshold and is properly proximal.
\end{itemize}

\medskip\noindent\textbf{Step 2: Inductive Amalgamated Decomposition}

Assume the theorem holds for all irreducible graphs with $k$ vertices, where $3 \le k < n$. Let $\Gamma$ be an irreducible graph with $n \ge 4$ vertices. Since $\Gamma^c$ is a connected graph with at least 4 vertices, it must contain a non-cut vertex $v$. Thus, the induced subgraph $\Gamma' = \Gamma \setminus \{v\}$ remains irreducible and has $n-1 \ge 3$ vertices. By the induction hypothesis, the subalgebra $A = M_{\Gamma'}$ is properly proximal.

Let $L = \operatorname{lk}_\Gamma(v)$ denote the link of $v$ in $\Gamma$. In the complement graph $G = \Gamma^c$, the vertices of $L$ are exactly the non-neighbors of $v$. Because $v$ must have at least one neighbor in the connected graph $G$, the set $L$ is strictly smaller than $\VV(\Gamma')$.
The full graph product $M_\Gamma$ natively decomposes into the amalgamated free product:
\[
M_\Gamma = A \ast_C B, \quad \text{where } C = M_L \text{ and } B = C \overline{\otimes} M_v.
\]

\medskip\noindent\textbf{Step 3: Geometric Ping-Pong on $\mathbb{B}(L^2M_\Gamma)$}

By Boutonnet, Ioana, and Peterson (2021), a tracial von Neumann algebra $M$ is properly proximal if there does not exist an $M$-central state $\varphi$ on $\mathbb{B}(L^2M)$ extending the trace $\tau_M$ such that $\varphi(T) = 0$ for all operators $T \in \mathcal{J}_{\mathrm{prop}}(M)$, where $\mathcal{J}_{\mathrm{prop}}(M)$ is the norm-closed proper proximal ideal generated by $\{x e_D^M y \mid x, y \in M, D \subset M \text{ an amenable von Neumann subalgebra}\}$.

Assume for the sake of contradiction that $M_\Gamma$ is not properly proximal. Then there exists such an $M_\Gamma$-central state $\varphi$ on $\mathbb{B}(L^2M_\Gamma)$ that completely vanishes on $\mathcal{J}_{\mathrm{prop}}(M_\Gamma)$.

Let $L^2M_\Gamma$ decompose into the orthogonal direct sum of $L^2C$ and the spaces of reduced alternating words over the orthocomplements $A \ominus C$ and $B \ominus C$. Define $Q_A$ as the orthogonal projection onto the closed span of all alternating words whose leftmost letter belongs to $A \ominus C$. Define $Q_B$ as the orthogonal projection onto the span of all alternating words whose leftmost letter belongs to $B \ominus C$. The alternating word basis strictly enforces the exact orthogonal equality $1 = e_C + Q_A + Q_B$.

Let $u_v \in M_v$ be the given trace-zero unitary. Define $u = 1 \otimes u_v \in C \overline{\otimes} M_v = B$. Since $\tau(u_v)=0$, we have $E_C(u)=0$, ensuring $u \in B \ominus C$. Left multiplication of any alternating word starting with $A \ominus C$ by $u$ strictly increases the word length by 1, producing a valid alternating word that begins with $u \in B \ominus C$. Thus, $u \operatorname{ran}(Q_A) \subset \operatorname{ran}(Q_B)$. Since $u$ is unitary, this ensures the projection inequality $u Q_A u^* \le Q_B$. Applying the $M_\Gamma$-central state $\varphi$, we deduce $\varphi(Q_A) = \varphi(u Q_A u^*) \le \varphi(Q_B)$.

Because $v$ is a non-cut vertex of $G = \Gamma^c$, we select specific vertices in $G$ to construct orthogonal shifts. Let $N_{G}(v)$ be the neighbors of $v$ in $G$. Elements in $N_G(v)$ are connected to $v$ in $G$, meaning they are strictly disconnected from $v$ in $\Gamma$, so they do not belong to $L$ (hence $M_{N_G(v)} \not\subset C$). 
\begin{itemize}
    \item \textbf{Case A ($|N_G(v)| \ge 2$):} Choose distinct $x, y \in N_G(v)$. 
    
  If $x \sim y$ in $G$ ($x \nsim y$ in $\Gamma$): Set $w_1 = u_x$ and $w_2 = u_x u_y u_x$. Because $x \notin L$, $E_C(w_1) = 0$, so $w_1 \in A \ominus C$. In the graph product, any reduced word containing at least one letter outside $L$ is strictly orthogonal to $M_L$. Since $x \nsim y$ in $\Gamma$, $u_x$ and $u_y$ do not commute, meaning $u_x u_y u_x$ is a reduced word. Because it contains $x \notin L$, its conditional expectation onto $C$ is exactly $E_C(w_2) = 0$, so $w_2 \in A \ominus C$. Further, $w_2^* w_1 = u_x^* u_y^* u_x^* u_x = u_x^* u_y^*$. This is a reduced word with letters outside $L$, guaranteeing $E_C(w_2^* w_1) = 0$.
  
  If $x \nsim y$ in $G$ ($x \sim y$ in $\Gamma$): Set $w_1 = u_x$ and $w_2 = u_y$. Both are strictly in $A \ominus C$. Since they commute, $w_2^* w_1 = u_y^* u_x$. Because $x \neq y$ and both $\notin L$, this is a product of mutually commuting elements strictly outside $L$, guaranteeing $E_C(w_2^* w_1) = E_C(u_y^*)E_C(u_x) = 0$.
  
    \item \textbf{Case B ($|N_G(v)| = 1$):} Let $x$ be the unique neighbor. Since $|V(G)| \ge 4$, $x$ must have another neighbor $y \neq v$ in $G$. Because $x$ is the unique neighbor of $v$ in $G$, $y \notin N_G(v)$, meaning $y \in L$ (thus $M_y \subset C$). Set $w_1 = u_x$ and $w_2 = u_x u_y u_x$. Then $w_1 \in A \ominus C$. For $w_2$, since $x \sim y$ in $G$ ($x \nsim y$ in $\Gamma$), $u_x$ and $u_y$ do not commute. The word $u_x u_y u_x$ is strictly reduced in the graph product and contains $x \notin L$, enforcing $E_C(w_2) = 0$, so $w_2 \in A \ominus C$. Furthermore, $w_2^* w_1 = u_x^* u_y^*$. This is a reduced word containing $x \notin L$, guaranteeing $E_C(w_2^* w_1) = 0$.
\end{itemize}

In all cases, we rigorously secure unitaries $w_1, w_2 \in A \ominus C$ satisfying $E_C(w_2^* w_1) = 0$. 
For any $\xi, \zeta \in \operatorname{ran}(Q_B)$, they are spanned by alternating words whose first letter is in $B \ominus C$. Left-multiplying by $w_1 \in A \ominus C$ strictly increases the alternating word length by 1, mapping $\xi$ to an alternating word starting with $A \ominus C$. Thus, $w_1 Q_B w_1^* \le Q_A$. Moreover, orthogonality is strictly preserved: $\langle w_1 \xi, w_2 \zeta \rangle = \tau(\zeta^* w_2^* w_1 \xi) = 0$. This holds because $E_C(w_2^* w_1) = 0$, ensuring that the middle product $w_2^* w_1 \in A \ominus C$, so $\zeta^* (w_2^* w_1) \xi$ remains a valid alternating word of length $\ge 3$, which necessarily has trace zero.

Thus, $w_1 Q_B w_1^*$ and $w_2 Q_B w_2^*$ are mutually orthogonal subprojections of $Q_A$. 
Applying $\varphi$, we acquire $2 \varphi(Q_B) = \varphi(w_1 Q_B w_1^*) + \varphi(w_2 Q_B w_2^*) \le \varphi(Q_A)$. Combined with the established $\varphi(Q_A) \le \varphi(Q_B)$, this algebraically forces $\varphi(Q_A) = 0$ and $\varphi(Q_B) = 0$. 
Consequently, $\varphi(e_C) = \varphi(1 - Q_A - Q_B) = 1$. Because $C \subset A$, the projection geometry strictly demands $e_C \le e_A$, enforcing $\varphi(e_A) = 1$.

\medskip\noindent\textbf{Step 4: The Algebraic Lifting Contradiction}

Define the functional $\psi$ on $\mathbb{B}(L^2A)$ by $\psi(T) = \varphi(e_A T e_A)$, utilizing the exact spatial identification $\mathbb{B}(L^2A) \cong e_A \mathbb{B}(L^2M_\Gamma) e_A$. 
Because $\varphi(e_A) = 1$, $\psi$ is a valid state on $\mathbb{B}(L^2A)$. For any $a \in A$, $a$ commutes with $e_A$, ensuring $\psi(a T) = \varphi(e_A a T e_A) = \varphi(a e_A T e_A) = \varphi(e_A T e_A a) = \psi(T a)$. Thus, $\psi$ is strictly $A$-central. Moreover, since $\varphi(e_A) = 1$, we unconditionally have $\varphi(1 - e_A) = 0$. The Cauchy-Schwarz inequality for states ensures $|\varphi(a (1 - e_A))|^2 \le \varphi(a a^*) \varphi(1 - e_A) = 0$, rigorously guaranteeing $\psi(a) = \varphi(e_A a e_A) = \varphi(a e_A) = \varphi(a) = \tau_M(a) = \tau_A(a)$, meaning $\psi$ extends the standard trace on $A$. 

We now evaluate $\psi$ on the proper proximal ideal $\mathcal{J}_{\mathrm{prop}}(A)$. By definition, $\mathcal{J}_{\mathrm{prop}}(A)$ is the norm-closed ideal in $\mathbb{B}(L^2A)$ generated by operators of the form $X e_D^A Y$, where $X, Y \in A$ and $D \subset A$ is an amenable von Neumann subalgebra. Under the spatial identification, $e_D^A$ corresponds precisely to the restriction of $e_D^M$, so on $L^2M_\Gamma$ it acts strictly as $e_D^M e_A = e_A e_D^M e_A$. 
For any generator $X e_D^A Y$ with $X, Y \in A$, we evaluate:
\[
\psi(X e_D^A Y) = \varphi(e_A X e_D^M Y e_A).
\]
Since $X, Y \in A$ and $A$ is a von Neumann subalgebra, the orthogonal projection $e_A$ strictly commutes with $X$ and $Y$ (which inherently act by left multiplication on $L^2M_\Gamma$), allowing us to rewrite the argument as $X e_A e_D^M e_A Y = X e_D^M Y$. 
Thus, $\psi(X e_D^A Y) = \varphi(X e_D^M Y)$.

Because $D \subset A \subset M_\Gamma$ is an amenable subalgebra of $M_\Gamma$ and $X, Y \in A \subset M_\Gamma$, the operator $X e_D^M Y$ inherently belongs to the proper proximal ideal $\mathcal{J}_{\mathrm{prop}}(M_\Gamma)$ strictly by definition. 
By our founding contradiction assumption, $\varphi$ vanishes entirely on $\mathcal{J}_{\mathrm{prop}}(M_\Gamma)$, which rigorously dictates $\varphi(X e_D^M Y) = 0$. 
By linearity and norm-continuity, $\psi$ completely vanishes on all of $\mathcal{J}_{\mathrm{prop}}(A)$.

Thus, $\psi$ is an $A$-central state on $\mathbb{B}(L^2A)$ extending $\tau_A$ that definitively vanishes on $\mathcal{J}_{\mathrm{prop}}(A)$. However, by the induction hypothesis, $A$ is properly proximal, mathematically prohibiting the existence of exactly such a state. 
This is a direct and fatal contradiction. Therefore, no such state $\varphi$ can exist, definitively establishing that the graph product von Neumann algebra $M_\Gamma$ is properly proximal.
\end{proof}

\end{document}